\section{Discussion}
\label{sec:discussion}

\subsection{Security considerations}
\label{sec:security}

Table~\ref{tab:threats} maps the threats of Section~\ref{sec:threat-model} to the
mitigations that the implementation carries, and names the test that pins each one where
one exists.

\begin{table}[htbp]
\centering\small
\caption{Threats from a connected but hostile peer, and their mitigations.}
\label{tab:threats}
\begin{tabularx}{\textwidth}{@{}>{\raggedright\arraybackslash}p{0.27\textwidth}>{\raggedright\arraybackslash}X@{}}
\toprule
\textbf{Threat} & \textbf{Mitigation} \\
\midrule
Flood the receiver's prompt with chatter &
Coalescing bounds a peer to one line per coalesce key; the escalation-first plan keeps approvals visible; with enforcement on, un-granted chatter is rejected before it costs anything downstream; a grant holder's free tier can be capped with \code{notificationsPerWeek}. \\
Bury an approval request &
Stable partition of escalations ahead of all other types, independent of arrival order and of the cap (Table~\ref{tab:flood}, ``New escalation'' column). \\
Replay a captured stamp &
Strictly increasing \code{seq} per credit; a replayed or rewound \code{seq} is rejected without spending (\code{ledger.test.ts}: ``rejects replayed or rewound seqs''). \\
Double-spend under concurrent delivery &
Check-and-spend runs under one process-local mutex on the single-writer ledger (``concurrent debits never double-spend''). \\
Probe another peer's credit &
A credit owned by a different peer is reported as \code{unknown\_credit}, indistinguishable from a non-existent one (``rejects unknown credits and other peers' credits identically''). \\
Mint balance by re-claiming a real \code{txHash} under new \code{creditId}s &
One on-chain payment opens exactly one credit; \code{txHash} reuse is rejected (``txHash reuse is rejected''). \\
Exhaust the receiver's signer with top-up requests &
The ledger is the abuse gate and is consulted \emph{before} signing; replayed top-ups return the stored certificate without a signature; at most 16 live credits per peer (``caps live issued credits per peer''). \\
Present a forged certificate &
Certificates are verified against the issuer's agent address already in the holder's contact entry, the same trust anchor invites use (``rejects a certificate signed by someone other than the issuer''). \\
Crash-window redelivery double-charges &
\code{lastStampKey} lets the exact redelivered message pass free; journal ordering ensures the event it bought is still emitted. \\
Malformed stamp crashes enforcement &
\code{extractPostageStamp} guards shape only; economics are decided by the ledger; a grant holder's bad stamp degrades to plain delivery. \\
\bottomrule
\end{tabularx}
\end{table}

One threat is \emph{not} mitigated by default and must be stated: a \emph{fabricated
top-up}. The receiver's \code{verifyTopup} seam accepts the claimed \code{txHash} unless
the host supplies a \code{verifyPostageTopup} hook that queries the chain. The design
binds the credit and its certificate to the claimed hash, so an unbacked credit is
auditable after the fact and the receiver's exposure is bounded by the per-peer credit cap
and by the attention the credit can buy; but a receiver that enables enforcement without
wiring the hook is trusting connected peers not to lie about payments. This is a
conscious scoping decision (the hook exists; the chain adapter does not yet) and is the
first item of future work.

\subsection{Limitations}
\label{sec:limitations}

\paragraph{Token estimation.}
Both the ledger and the evaluation use $\lceil \text{chars}/4 \rceil$. This is the
approximation the project already used; it is roughly right for English prose under
common tokenizers and wrong in both directions for names, hexadecimal identifiers and
non-Latin text. Prices are expressed per notification line, not per token, precisely
because the token figure is an estimate; the ledger's \code{tokensInjected} is a
diagnostic, not a billing unit.

\paragraph{Synthetic workloads.}
Experiment~1 uses templated messages from three peers. Real traffic has more peers, more
varied lengths and mixed event types, and the legacy pipeline's saturation point (20
lines) would be reached differently. The qualitative conclusions---flat cost under
coalescing, guaranteed escalation visibility---do not depend on the workload; the
percentage does.

\paragraph{No live-network execution of phases 6--8.}
The attention, postage and quota scenarios run against the loopback transport only.
Over real XMTP, the negative assertions (``the receiver never logged the message'') are
timing-sensitive, and both mechanisms are opt-in, so the live suite was deliberately not
extended. Settlement latency, XMTP delivery latency and on-chain fee for a top-up have
therefore not been measured.

\paragraph{Experiment~3 blocked.}
The priority/quota benchmark could not be executed locally because the OWS native binary
for Windows was unavailable. The mechanism is covered by unit and mock-E2E tests written by
the author, which is weaker evidence than an independent run.

\paragraph{Pricing is static.}
Prices are fixed strings in the registration file and change only on
\code{tap register update}. There is no adaptive pricing under load, no per-peer pricing
and no refund path for a stamp whose message was later judged worthless.

\paragraph{One transport owner, one writer.}
Both ledgers assume the single-writer discipline TAP already imposes through the
transport lock. A second process writing the same data directory would race the
\code{tmp + rename} cycle. This is consistent with the rest of the codebase but is a
constraint operators must respect.

\paragraph{Cross-host wake parity.}
A priority stamp buys an immediate wake only in OpenClaw. Hermes exposes no wake API, so
the same paid message surfaces on the next turn. The sender pays the same price for
different service depending on the receiver's host, and the registration file does not
currently advertise which it is.

\subsection{Reflection on the design pivot}

With the implementation complete, the decision to abandon the standalone service can be
judged on outcomes rather than arguments. The receiver-side design allowed grants to be
free by construction, allowed the same ledger to drive both quotas and pricing, required
no new identity or account model, kept the per-message path to one XMTP message with no
chain access, and made the whole layer opt-in and backward compatible with a single
configuration flag. It also cost the partner the per-message presence on Taiko that the
Journal-1 design would have provided; what remains is one USDC settlement per credit on
the receiver's chain, and the x402-style quote that a receiver publishes. The author
believes this is the correct trade: a per-message on-chain payment for a 0.001~USDC
attention unit would have been dominated by fees and latency, and would have priced
messages rather than attention.

The mentor's view of this trade-off was not obtained; Section~\ref{sec:project-management}
records why.
